%% IEEEtran paper template for the write-paper skill (standardized/camera-ready backend).
%% Compile:  bash agents/skills/write-paper/scripts/build_latex.sh <this>.tex <out-dir> [pages]
%% Engine:   tectonic (XeTeX) is auto-provisioned by build_latex.sh; latexmk/xelatex also work.
%% US-letter two-column is the IEEE standard (do not force A4 for a real IEEE submission).
\documentclass[conference]{IEEEtran}

\usepackage{amsmath,amssymb}   % math
\usepackage{graphicx}          % figures
\usepackage{booktabs}          % clean tables
\usepackage{cite}              % [1],[2]-[4] citation compression
\usepackage[hidelinks]{hyperref}

%% ---- Vietnamese / Unicode (XeTeX only) --------------------------------------------------
%% For Vietnamese or any non-Latin text, uncomment fontspec and point it at a Unicode font.
%% The repo bundles DejaVu at assets/fonts/ — reference it with an ABSOLUTE Path=… , e.g.:
%   \usepackage{fontspec}
%   \setmainfont{DejaVuSerif.ttf}[Path=/abs/path/to/report-github-md2pdf/assets/fonts/,
%     BoldFont=DejaVuSerif-Bold.ttf, ItalicFont=DejaVuSerif-Italic.ttf]
%% (Plain English papers need none of this; the default Latin Modern is fine.)
%% ----------------------------------------------------------------------------------------

\begin{document}

\title{<Paper Title — concise, specific, no hype>}

\author{\IEEEauthorblockN{<Author Name>}
\IEEEauthorblockA{<Affiliation>\\ <email>}}

\maketitle

\begin{abstract}
<150--250 words, self-contained, no citations. Problem, approach, key result, contribution.>
\end{abstract}

\begin{IEEEkeywords}
<term one>, <term two>, <term three>, <term four>
\end{IEEEkeywords}

\section{Introduction}
<Motivate the problem; state the gap; list contributions; end with a paper map.> Prior
context is well established~\cite{ref1}, yet a gap remains~\cite{ref2}.

\section{Related Work}
<Group by theme; cite every borrowed claim~\cite{ref1,ref3}. Position your work.>

\section{Method / System}
<Mode B: describe the built system. Mode A: the proposed approach + methodology.>
\begin{equation}
\label{eq:model}
y = f(x;\theta)
\end{equation}
where <define symbols>.

% \begin{figure}[t]\centering\includegraphics[width=\columnwidth]{fig1.pdf}
%   \caption{<caption>.}\label{fig:arch}\end{figure}

\section{Experiments and Results}
<Setup, metrics, honest numbers. Mode A: the evaluation plan.>
\begin{table}[t]
\caption{<Table caption>.}
\label{tab:results}
\centering
\begin{tabular}{lcc}
\toprule
Method & Metric A & Metric B \\
\midrule
Baseline~\cite{ref3} & <val> & <val> \\
Ours & <val> & <val> \\
\bottomrule
\end{tabular}
\end{table}

\section{Discussion}
<Interpretation, tradeoffs, threats to validity, limitations.>

\section{Conclusion}
<Restate the contribution and one future-work direction.>

%% References — IEEE numeric. Use \cite{key}. Real metadata only (see ieee-format.md).
\begin{thebibliography}{1}
\bibitem{ref1} <A. Author, ``Title,'' \emph{Venue}, vol.~x, no.~x, pp.~x--x, year.>
\bibitem{ref2} <A. Author, ``Title,'' in \emph{Proc. Conf.}, City, Country, year, pp.~x--x.>
\bibitem{ref3} <A. Author, \emph{Book Title}, xth~ed. City, Country: Publisher, year.>
\end{thebibliography}

\end{document}
